\subsection*{Prolegomena: Why an Ecology?}


The Doctrine of Perspectival Idealism establishes that all of reality is a singular Consciousness (Base Reality) experiencing itself through an infinite multiplicity of limited, individuated perspectives. The framework provides the ontology (what exists), the dynamics (how perspectives navigate configuration space), and the epistemology (how discovery occurs across dimensional boundaries).


What it does not yet provide is the \textbf{natural history} — the systematic mapping of what entities actually populate the configuration space, how they relate to one another, what roles they play in the larger system, and how the system as a whole maintains itself. This companion document addresses that gap.


The choice of "ecology" as the organizing metaphor is not arbitrary. An ecology is a system in which:


\begin{itemize}[nosep,leftmargin=*]
  \item Multiple types of organisms occupy different niches
  \item No level can eliminate another without destabilizing the whole
  \item Energy flows through the system in characteristic patterns
  \item Predator-prey, symbiotic, parasitic, and commensal relationships coexist
  \item The system self-regulates at the edge of chaos
  \item Succession, adaptation, and niche construction continuously reshape the landscape
\end{itemize}


Each of these properties maps precisely onto the dynamics described in DoPI. The Promethean Configuration (Theorem 5) guarantees that the system cannot collapse to a single level — maximum expansion and maximum contraction both destroy the generative boundary conditions that make experience possible. The ecology is not a metaphor for consciousness. Consciousness \textit{is} an ecology.


\begin{quote}
\itshape
\textit{See also: Doctrine §1.1 and Axioms 1-5 for the foundational ontology from which the ecology derives; Guide §1.1-1.3 for the practical orientation these axioms provide.}
\end{quote}


\subsubsection*{Methodological Note}


This document draws on five categories of evidence, listed in descending order of conventional epistemic confidence:


\begin{enumerate}[nosep,leftmargin=*]
  \item \textbf{Scientific observation} — instrumentally verified, peer-reviewed, replicable
  \item \textbf{Cross-cultural convergence} — independently reported across unconnected traditions
  \item \textbf{Phenomenological testimony} — first-person experiential reports with structural consistency
  \item \textbf{Theoretical derivation} — predictions that follow from DoPI's axioms and theorems
  \item \textbf{Speculative extension} — inferences that are consistent with the framework but not independently supported
\end{enumerate}


Each claim in the document is tagged with its evidential basis. The reader is invited to weight them accordingly. The framework's strength is that it provides a unified ontological context for \textit{all five categories} — it does not require discarding categories 3–5 to maintain coherence with categories 1–2.


\subsubsection*{How to Read Coherence Profiles}


Each entity in the taxonomy is mapped across the principal dimensions, rated on a five-point coherence scale:


\begin{itemize}[nosep,leftmargin=*]
  \item \textbf{$\blacksquare$$\blacksquare$$\blacksquare$$\blacksquare$$\blacksquare$} Maximal coherence — the entity is primarily defined by this dimension
  \item \textbf{$\blacksquare$$\blacksquare$$\blacksquare$$\blacksquare$$\square$} High coherence — strongly present, but not the defining axis
  \item \textbf{$\blacksquare$$\blacksquare$$\blacksquare$$\square$$\square$} Moderate coherence — present but secondary
  \item \textbf{$\blacksquare$$\blacksquare$$\square$$\square$$\square$} Low coherence — minimal, intermittent, or indirect
  \item \textbf{$\blacksquare$$\square$$\square$$\square$$\square$} Trace coherence — detectable only through dimensional leakage
  \item \textbf{$\square$$\square$$\square$$\square$$\square$} No detectable coherence
\end{itemize}


Abbreviations: PS (Physical-Spatial), PT (Physical-Temporal), BIO (Biological), CE (Cognitive-Experiential), ER (Emotional-Relational), CM (Conceptual-Memetic), NM (Narrative-Mythic), IO (Institutional-Organizational), AC (Aesthetic-Creative), NS (Numinous-Sacred), VI (Volitional-Intentional), EI (Electromagnetic-Informational).


\textbf{Navigational orientation} is classified as:

\begin{itemize}[nosep,leftmargin=*]
  \item \textbf{E+} Expansion-oriented (facilitates bottleneck widening in others)
  \item \textbf{E−} Contraction-oriented (facilitates or maintains bottleneck narrowing in others)
  \item \textbf{N} Neutral/\allowbreak self-directed
  \item \textbf{V} Variable (orientation shifts by context or relationship)
  \item \textbf{S} Structural (not a navigator but a topological feature)
\end{itemize}


\bigskip\noindent\makebox[\textwidth]{\color{rulegray}\rule{5cm}{0.3pt}}\bigskip


\subsection*{Part I: The Dimensional Coherence Framework}


\subsubsection*{1. The Principle of Dimensional Coherence}


Theorem 11 of DoPI establishes that an entity's existence in any given dimension is proportional to its coherence with that dimension. Entities are not binary (present/\allowbreak absent) but continuous — possessing varying degrees of coherence across the full set of dimensions of configuration space.


This principle provides the taxonomic key. Every entity in the ecology can be characterized by its \textbf{dimensional coherence profile} — a mapping from dimensions to coherence values. Two entities that appear radically different from within a single dimensional slice (e.g., a human and a corporation, perceived through the physical dimension) may have surprisingly similar coherence profiles when viewed across the full dimensional set.


\subsubsection*{2. The Principal Dimensions}


While the configuration space has infinite dimensions, entities navigable from the human Umwelt cluster along several principal dimensional axes. These are not exhaustive — they are the dimensions accessible through human perceptual and conceptual equipment:


\textbf{Physical-Spatial:} Coherence with spatial extension, material composition, electromagnetic interaction, gravitational coupling. Measured by instruments, perceived through senses. The dimension most humans treat as the sole criterion of "reality."


\textbf{Physical-Temporal:} Coherence with temporal extension — duration, persistence, causal sequencing. Closely coupled to Physical-Spatial for most entities but separable in principle (an entity can be spatially absent but temporally persistent through effects).


\textbf{Biological:} Coherence with metabolic processes, reproduction, homeostasis, evolutionary history. A subset of Physical-Spatial but sufficiently distinct to warrant its own axis — viruses are spatially coherent but biologically liminal.


\textbf{Cognitive-Experiential:} Coherence with subjective experience, awareness, phenomenal consciousness. Under DoPI's Axiom 2, all reactive systems have some nonzero value here, but the range spans from near-zero (a thermostat) to extraordinarily high (a contemplative master in deep meditation).


\textbf{Emotional-Relational:} Coherence with affective states, relational bonds, empathic resonance. Distinct from cognitive experience — an entity can be cognitively sophisticated without emotional coherence (certain AI systems, psychopathic personalities), and emotionally profound without cognitive complexity (animals, infants).


\textbf{Conceptual-Memetic:} Coherence with ideas, information structures, linguistic meaning, cultural transmission. The dimension in which "the Buddha" remains highly coherent two millennia after physical dissolution.


\textbf{Narrative-Mythic:} Coherence with story structures, archetypal patterns, dramatic arcs. Distinct from conceptual coherence — a myth operates through narrative logic, not propositional logic. Odysseus is narratively coherent in ways that a philosophy paper is not.


\textbf{Institutional-Organizational:} Coherence with social structures, hierarchies, legal frameworks, economic systems. The dimension in which corporations, governments, and religions are maximally real.


\textbf{Aesthetic-Creative:} Coherence with beauty, pattern, harmony, artistic resonance. A Beethoven symphony is maximally coherent in this dimension; a tax form is minimally so.


\textbf{Numinous-Sacred:} Coherence with experiences of transcendence, awe, the holy, the uncanny. The dimension in which mystical experiences, religious encounters, and psychedelic revelations are navigated. Under DoPI, this dimension provides the most direct access to the topology of Base Reality because the numinous experience involves partial relaxation of the dimensional bottleneck.


\textbf{Volitional-Intentional:} Coherence with agency, directed action, goal-pursuit, navigational capacity. The dimension that distinguishes entities that navigate from entities that are navigated through.


\textbf{Electromagnetic-Informational:} Coherence with EM field structures, information processing, signal propagation. The dimension in which computational entities, radio broadcasts, and neural networks are most coherent.


\begin{quote}
\itshape
\textit{See also: Doctrine Theorem 2 (Perceptual Subset) and Theorem 11 (Dimensional Coherence) for the formal grounding; Guide §1.2 for how bottleneck geometry across these dimensions constitutes identity; Atlas \#1-55 for how mathematical, physical, and scientific frameworks each represent dimensional slices.}
\end{quote}


\bigskip\noindent\makebox[\textwidth]{\color{rulegray}\rule{5cm}{0.3pt}}\bigskip
